% ============================================================================
%  CAPÍTULO IV: CONCLUSIONES
% ============================================================================

\section{Conclusiones}

La implementación del servidor web permitió comprender de manera integral el
funcionamiento de este servicio fundamental en una infraestructura de red. A
continuación se presentan las conclusiones más relevantes.

\subsection{Resultados obtenidos}

El servidor web fue implementado exitosamente con las siguientes
características verificadas:

\begin{itemize}
  \item Nginx operativo como servidor de contenido estático y proxy inverso,
        atendiendo solicitudes HTTP y HTTPS en los puertos 80 y 443.
  \item Certificado SSL/TLS auto-firmado generado automáticamente al
        iniciar el contenedor, cubriendo el dominio base y todos los
        subdominios del proyecto mediante \textit{wildcard}.
  \item Redirección HTTP a HTTPS funcional en todos los virtual hosts,
        garantizando que las comunicaciones se cifren.
  \item Cinco proxies inversos configurados y operativos: CUPS (impresión),
        Filebrowser (archivos), Portainer (gestión Docker), Zabbix (monitoreo)
        y el sitio principal.
  \item Soporte WebSocket para Portainer con timeouts extendidos.
  \item Resolución dinámica de upstreams mediante el DNS embebido de Docker,
        garantizando la resiliencia del servidor ante reinicios de backends.
  \item Servidor por defecto (\texttt{default\_server}) que captura
        peticiones por IP o dominios no declarados.
\end{itemize}

\subsection{Dificultades encontradas}

Durante la configuración se encontraron las siguientes dificultades:

\begin{itemize}
  \item Certificados auto-firmados: el uso de certificados SSL
        auto-firmados genera avisos de seguridad en los navegadores web,
        requiriendo la aceptación manual de excepciones. En un entorno de
        producción se recomienda utilizar certificados emitidos por una
        autoridad de certificación confiable.
  \item Resolución de upstreams en el arranque: Nginx falla si no puede
        resolver un nombre de upstream al iniciar. Se solucionó utilizando
        variables y el resolver de Docker (\texttt{127.0.0.11}), que permite
        la resolución en tiempo de ejecución.
  \item Soporte WebSocket para Portainer: la interfaz de Portainer
        requiere conexiones WebSocket persistentes. Sin la configuración de
        las cabeceras \texttt{Upgrade} y \texttt{Connection}, la interfaz no
        funciona correctamente.
  \item Tamaño de subida en Filebrowser: el límite predeterminado de
        Nginx para el cuerpo de las solicitudes (1\,MB) era insuficiente para
        la subida de archivos. Se solucionó con
        \texttt{client\_max\_body\_size 500M}.
\end{itemize}

\subsection{Importancia del servidor web en la administración de redes}

El servidor web aporta los siguientes beneficios a la administración de
sistemas:

\begin{itemize}
  \item Punto de entrada único: como proxy inverso, centraliza el acceso
        a todos los servicios del proyecto bajo un único dominio con múltiples
        subdominios.
  \item Seguridad: la terminación TLS en Nginx cifra todas las
        comunicaciones entre los clientes y los servicios internos, incluso
        cuando los backends no implementan cifrado propio.
  \item Flexibilidad: los virtual hosts permiten agregar nuevos servicios
        sin modificar la infraestructura de red ni abrir puertos adicionales
        en el firewall.
  \item Presencia institucional: el sitio web estático del proyecto
        proporciona una interfaz profesional accesible desde cualquier
        navegador.
  \item Simplicidad operativa: al concentrar la configuración TLS y el
        enrutamiento en un solo punto, se simplifica la administración de
        certificados y la gestión del acceso a los servicios.
\end{itemize}

\subsection{Aprendizajes adquiridos}

El taller permitió desarrollar competencias prácticas en administración de
servicios de red:

\begin{itemize}
  \item Se comprendió el rol del servidor web como proxy inverso y
        terminador TLS, consolidando el acceso a múltiples servicios bajo una
        única interfaz HTTPS.
  \item Se aprendió a configurar virtual hosts en Nginx para servir
        múltiples subdominios con distintos backends.
  \item Se implementó la generación automática de certificados SSL
        auto-firmados, comprendiendo los conceptos de clave pública,
        \texttt{subjectAltName} y cadenas de confianza.
  \item Se comprendió la importancia de la resolución dinámica de upstreams
        en entornos contenerizados, donde las direcciones IP de los servicios
        pueden cambiar entre reinicios.
  \item Se aprendió a configurar soporte WebSocket en Nginx para servicios
        que requieren comunicación bidireccional en tiempo real.
  \item Se valoró la infraestructura como código: toda la configuración
        del servidor web quedó definida en archivos versionados,
        permitiendo su reproducción exacta en cualquier entorno.
\end{itemize}
